\scanpage{37}
squelette de dim 1'' (1-complexe combinatoire, muni du système des trois \deleted{ps.} sous 1-complexes $D_i$. \deleted{On a} Donc $K$ est la réunion des sous\deleted{\unsure{-}} 1-complexes $D_i$, qui sont des bigones, \inserted{chacun} définis par des ens \inserted{à deux él.} de sommets \inserted{$S_i$} et d'arêtes \inserted{$A_i$} \deleted{donnés} \deleted{par} \inserted{avec}
\[
 \deleted{\{S_i\simeq I\setminus\{i\},\ A_i}
\]
\[
 S_i\simeq I-\{i\}
\]
(NB. \deleted{\unsure{Car}} l'ens des sommets de $K$ (i.e. de $\unsure{\mathcal C|C}$) est en corr. biun. avec $I$, en associant à chaque $i\in I$ l'unique sommet qui n'est pas dans $D_i$). Soit
\[
 \Omega_i=S_i\wedge A_i=(I-\{i\})\wedge A_i
\]
l'ens des orientations du bigône $D_i$, i.e. de la ps. droite $\deleted{\unsure{A}}D_i$. Soit
\[
\begin{aligned}
 \Phi_i&=\Omega_i\wedge\bigwedge_{\alpha\in I\setminus\{i\}}(I\setminus\{\alpha\})\\
 &\simeq A_i\wedge\deleted{\unsure{(}}\bigwedge_{\alpha\in I}(I\setminus\{\alpha\})
\end{aligned}
\]
de sorte que l'on a des isom canoniques
\[
 \bigwedge_{i\in I}\Omega_i\simeq\bigwedge_{i\in I}\Phi_i
 \simeq\left(\bigwedge_{i\in I}A_i\right)\wedge
 \underbrace{\bigwedge_{\alpha\in I}(I\setminus\{\alpha\})}_{\operatorname{or} I}\tag*{(15)}
\]
les ens $S_i,\deleted{\unsure{D}}A_i,\Omega_i,\Phi_i$ $(i\in I)$ et les isom
% archive sheet number: 36
